%% multicols互換の環境：\flushcolumns，\raggedcolumnsと，\nobreakの後で始める段組
%% ・表題の後に\vskipと\nobreakを置いてからmulticols環境を始める（前のアキがページに
%%   移っているので，そのアキを段組の前のアキに数える）
%% ・\flushcolumnsで始め，\raggedcolumnsで終える（満ページは伸ばし，揃えるページは
%%   伸ばさない）．行送りに伸びのある本文
\documentclass[fontsize=10pt,twocolumn]{jlreq}
\usepackage[debug,columns=1]{evenend}
\setlength\columnseprule{0.4pt}
\newcount\n
\newcommand\items[1]{\n=0 \loop\ifnum\n<#1 \advance\n1 項目\the\n\dotfill\the\n\par\repeat}
\begin{document}
\begingroup\centering\Large 表題\par\endgroup
\vskip 20pt\nobreak
\begingroup
\baselineskip 15pt plus 1pt
\parindent0pt
\flushcolumns
\begin{multicols}{3}
\items{230}
\raggedcolumns
\end{multicols}
\endgroup
後の段落．
\end{document}
